\title[Does Motivation Matter?]{Does Motivation Matter? A Systematic Review and Meta-Analysis of Outcomes Following Intentional Ingestion of Foreign Objects}

\author*[1,3,4]{\fnm{Jack} \sur{Galbraith-Edge}}\email{research@jmge.io}
\author[1,2]{\fnm{Giles} \sur{Cattermole}}\email{Giles.Cattermole@qmul.ac.uk}

\affil*[1]{\orgdiv{Blizzard Institute, Faculty of Medicine and Dentistry}, \orgname{Queem Mary, University of London}, \orgaddress{\street{4 Newark Street}, \city{London}, \postcode{E1 2AT}, \country{United Kingdom}}}

\affil[2]{\orgdiv{Emergency Department (PRUH)}, \orgname{King's College Hospital NHS Foundation Trust}, \orgaddress{\street{Denmark Hill}, \city{London}, \postcode{SE5 9RS}, \country{United Kingdom}}}

\affil[3]{\orgname{East of England Ambulance Service NHS Trust}, \orgaddress{\street{Whiting Way}, \city{Melbourn}, \state{Cambridgeshire}, \postcode{SG8 6EN}, \country{United Kingdom}}}

\affil[4]{\orgname{South Western Ambulance Service NHS Foundation Trust}, \orgaddress{\street{Eagle Way}, \city{Exeter}, \state{Devon}, \postcode{EX2 7HY}, \country{United Kingdom}}}

\abstract{\textbf{Background:} Intentional Ingestion of Foreign Object(s) (IIFO) is a distinct, increasingly recognised form of self-harm that poses complex management challenges. The effect of motivation on outcomes following IIFO has not been studied.

\textbf{Aims:} To synthesise evidence on outcomes following IIFO and determine whether patient motivation, object characteristics, or demographics influence the need for intervention, complications, or mortality. 

\textbf{Methods:} A systematic literature search identified studies reporting intentional ingestion of non-nutritive foreign bodies in humans. Eligible studies included any design and age group reporting endoscopic, surgical, or conservative outcomes. Exclusions were accidental or substance ingestion, non-English full texts, animal studies, and reports lacking motivation, object, or outcome data. Screening was performed by a single reviewer with 10\% double-checked. Data were extracted at both study and case level. Meta-analysis of proportions summarised outcomes, and associations were explored using $\chi^2$ testing and univariate meta-regression. 

\textbf{Results:} Seventy-one case reports ($n=71$) and three case series ($n=90$) were included. Pooled outcome rates were: endoscopy 50\%, surgery 30\%, conservative management 20\%, complications 30\%, and mortality $<1$\%.  Case series populations were uniformly male and detained, with protest motivations ($\approx80$\%) and sharp objects ($\approx90$\%) predominant. Protest motivation significantly reduced surgical intervention (OR = 0.98, 95\% CI 0.98--0.98, $p=0.003$). Intent-to-harm increased surgical risk six-fold (OR = 5.68, 95\% CI 1.43--22.64, $p=0.020$). Adults aged 40--64 had lower surgical rates but higher mortality, often with psychiatric comorbidity.

\textbf{Conclusion:} Motivation significantly influences IIFO outcomes. Protest ingestion may be managed conservatively, whereas intent-to-harm and psychiatric comorbidity indicate higher surgical risk. Improved prospective reporting is needed.}

\keywords{Self-harm, Foreign Body, Foreign Object, Ingestion, Motivation}
\maketitle

